\documentclass[11pt,a4paper]{article}

\usepackage{amsmath,amssymb,amsthm}
\usepackage{algorithm}
\usepackage{algpseudocode}
\usepackage{listings}
\usepackage{hyperref}
\usepackage{booktabs}
\usepackage{geometry}
\usepackage{xcolor}

\geometry{margin=2.5cm}

\lstset{
  language=C++,
  basicstyle=\ttfamily\small,
  keywordstyle=\bfseries\color{blue!70!black},
  commentstyle=\itshape\color{green!50!black},
  numbers=left,
  numberstyle=\tiny\color{gray},
  frame=single,
  breaklines=true,
  columns=flexible
}

\newtheorem{definition}{Definition}
\newtheorem{lemma}{Lemma}
\newtheorem{theorem}{Theorem}
\newtheorem{remark}{Remark}

\title{%
  Extended-Double-Double Floating-Point Arithmetic: \\
  A binary80-Based Two-Limb Expansion in the QD Library%
}
\author{%
  Nakata Maho \\
  \textit{RIKEN} \\
  \texttt{maho@riken.jp}%
}
\date{2026}

\begin{document}
\maketitle

\begin{abstract}
We describe the design and algorithmic details of \texttt{edd\_real},
an \emph{extended-double-double} floating-point type that
complements the \texttt{dd\_real}, \texttt{td\_real}, and
\texttt{qd\_real} families in the QD library by replacing the
IEEE~754 binary64 limb with the GNU C++ binary80 type
\texttt{\_Float64x}.  An \texttt{edd\_real} value is represented as
an unevaluated sum of two non-overlapping binary80 numbers
$a = x_0 + x_1$, yielding approximately 126-bit ($\approx 38$
decimal digit) precision while retaining the native binary80
dynamic range of $\approx 10^{\pm 4932}$.  Because the limb
mantissa width is $p = 64$ rather than $53$, the Dekker splitter
becomes $2^{32}+1$ and the multiplication, division, and square-root
kernels must scale by $2^{\pm 33}$ (or larger) near the top of the
exponent range to avoid overflow.  We give algorithms for splitting
with rescaling, two-component renormalization, addition,
subtraction, multiplication, overflow-safe division, Karp-style
square root, and a hybrid transcendental strategy that delegates
to \texttt{qd\_real} whenever the argument is within its safe range
and falls back to native Taylor/Newton schemes otherwise.
Interoperability with \texttt{dd\_real} and \texttt{qd\_real} is
provided by explicit renormalizing truncations in both directions.
The implementation is available as an optional component of the QD
library, enabled at configure time on GNU C++ x86 targets.
\end{abstract}

\section{Introduction}
\label{sec:intro}

Many applications in computational science require more than the
$\approx 16$ decimal digits offered by IEEE~754 binary64
(double-precision).  The QD library of Hida, Li, and
Bailey~\cite{QD,HidaLiBailey2000} responds to this need by
representing higher-precision values as unevaluated sums of
binary64 limbs: two limbs for \texttt{dd\_real}
($\approx 106$ bits, $\approx 32$ digits), four limbs for
\texttt{qd\_real} ($\approx 212$ bits, $\approx 63$ digits), and
three limbs for the intermediate \texttt{td\_real}
($\approx 159$ bits, $\approx 47$ digits).  All three rely
exclusively on binary64 arithmetic and error-free transformations
(\textsc{TwoSum}, \textsc{TwoProd}) and are portable across every
IEEE~754 platform.

Several GNU C++ targets (notably x86 and x86\_64) additionally
expose the IA-32 80-bit extended-precision format as the
\texttt{\_Float64x} type.  This format carries a 64-bit explicit
mantissa and an exponent range of $\pm 16\,383$, giving a single
limb about 19 decimal digits of precision and a dynamic range of
$\approx 10^{\pm 4\,932}$.  A two-limb expansion of
\texttt{\_Float64x} therefore reaches about 126 bits, or
$\approx 38$ decimal digits, occupying the gap between
\texttt{dd\_real} (32 digits) and \texttt{td\_real} (47 digits)
while simultaneously inheriting a far wider exponent range than
any of the binary64-based types.  We call the resulting type
\texttt{edd\_real} (\emph{extended double-double}).

Why is \texttt{edd\_real} interesting?  From a pure-precision
standpoint it delivers fewer decimal digits than
\texttt{td\_real}, so for applications that are dominated by the
length of the significand, \texttt{td\_real} remains preferable.
The practical value of \texttt{edd\_real} lies elsewhere:

\begin{itemize}
  \item \textbf{Arithmetic cost.}  \texttt{edd\_real} stores only
    two limbs, so its addition and multiplication cost follow the
    well-known \texttt{dd\_real} formulas (2--3
    \textsc{TwoSum}/\textsc{TwoProd} calls plus renormalization)
    rather than the larger-arity kernels of \texttt{td\_real} and
    \texttt{qd\_real}.  Each limb operation is a single binary80
    FPU instruction on x86.
  \item \textbf{Dynamic range.}  The binary80 base gives a
    representable range up to approximately $2^{16383} \approx
    10^{4\,932}$, versus $\approx 10^{308}$ for the binary64-based
    types.  This matters for applications that must keep both
    precision and very large or very small magnitudes, such as
    AGM-based computations of transcendental constants.
  \item \textbf{Drop-in interoperability.}  \texttt{edd\_real}
    shares the QD library's conventions (operator overloading,
    \texttt{sqrt}, \texttt{exp}, \texttt{log},~$\ldots$) and adds
    explicit conversion routines to and from \texttt{dd\_real}
    and \texttt{qd\_real}.
\end{itemize}

Naively porting \texttt{dd\_real} to binary80 is subtle.  Because
the mantissa now has $p = 64$ bits, the Dekker splitter is
$C = 2^{\lceil p/2 \rceil} + 1 = 2^{32} + 1$ rather than the
familiar $2^{27}+1$ used by \texttt{dd\_real}.  The product $C \cdot a$
can overflow when $|a|$ approaches the binary80 maximum, and the
same phenomenon reappears in division (where $a_{\mathrm{hi}} / b$
can overflow) and in square root (where $\mathrm{sqr}(s)$ can
overflow for large $s$).  Each kernel therefore carries an explicit
rescale-and-restore pattern, keyed to a different threshold.
Section~\ref{sec:split} makes these thresholds precise.

The remainder of the paper is organized as follows.
Section~\ref{sec:prelim} recalls the binary80 format and the
error-free transformations we rely on.
Section~\ref{sec:repr} fixes the representation and the
non-overlapping invariant.
Section~\ref{sec:split} presents the splitter with rescaling.
Sections~\ref{sec:renorm}--\ref{sec:sqrt} cover renormalization,
addition, subtraction, multiplication, division, and square
root.  Section~\ref{sec:trans} describes the hybrid transcendental
strategy.  Section~\ref{sec:mixed} gives conversions and
mixed-precision arithmetic.
Sections~\ref{sec:const}--\ref{sec:test} describe constants,
build-system detection, and the test suite.

\section{Preliminaries}
\label{sec:prelim}

\subsection{The binary80 format}

The Intel 80-bit extended-precision format, exposed by GNU~C++
as \texttt{\_Float64x}, is the reference format for all
\texttt{edd\_real} arithmetic.  Its parameters are:

\begin{center}
\begin{tabular}{lc}
\toprule
Parameter & Value \\
\midrule
Explicit mantissa bits $p$ & $64$ \\
Exponent bits & $15$ \\
Bias & $16\,383$ \\
$E_{\min},\ E_{\max}$ (macro) & $-16\,381,\ 16\,384$ \\
Machine epsilon $\varepsilon_0 = 2^{-63}$ & $\approx 1.08\times 10^{-19}$ \\
Decimal digits per limb & $\approx 19$ \\
\bottomrule
\end{tabular}
\end{center}

Following GNU C++ conventions, the \texttt{edd} backend refers to
these parameters as the preprocessor macros
\texttt{\_\_FLT64X\_MANT\_DIG\_\_ = 64} and
\texttt{\_\_FLT64X\_MAX\_EXP\_\_ = 16384}, which are verified at
\texttt{configure} time.  The builtins
\verb|__builtin_sqrtf64x|, \verb|__builtin_ldexpf64x|,
\verb|__builtin_fabsf64x|, \verb|__builtin_floorf64x|,
\verb|__builtin_copysignf64x|, and \verb|__builtin_frexpf64x|
are used throughout.  We introduce the abbreviation
\verb|edd_word| for \texttt{\_Float64x}.

\subsection{Error-free transformations}

All \texttt{edd\_real} kernels are built from the following
standard error-free transformations, reproduced here with the
exact operation counts used by the implementation.

\medskip

\noindent\textsc{TwoSum}$(a,b)$: for IEEE-rounded addition,
returns $(s, e)$ with $s = \mathrm{fl}(a+b)$ and $a+b = s+e$
exactly, whenever $s$ is finite.  Cost: 6 flops.
\begin{algorithmic}[1]
\State $s \gets a + b$
\State $b' \gets s - a$
\State $e \gets (a - (s - b')) + (b - b')$
\State \Return $(s, e)$
\end{algorithmic}

\noindent\textsc{QuickTwoSum}$(a,b)$: same guarantee when $|a| \ge
|b|$.  Cost: 3 flops.
\begin{algorithmic}[1]
\State $s \gets a + b$;\quad $e \gets b - (s - a)$
\State \Return $(s, e)$
\end{algorithmic}

\noindent\textsc{TwoDiff}$(a,b)$: analogous to \textsc{TwoSum} for
subtraction.

\medskip

\noindent\textsc{TwoProd}$(a,b)$: returns $(p, e)$ with $p =
\mathrm{fl}(ab)$ and $ab = p+e$ exactly.  On a target with a fused
multiply--add this is a single multiplication and a single FMA.
Without FMA, the Dekker--Veltkamp algorithm is used:
\begin{algorithmic}[1]
\State $p \gets a \cdot b$
\State $(a_h, a_l) \gets \textsc{Split}(a)$;\quad $(b_h, b_l) \gets \textsc{Split}(b)$
\State $e \gets ((a_h b_h - p) + a_h b_l + a_l b_h) + a_l b_l$
\State \Return $(p, e)$
\end{algorithmic}
The splitter \textsc{Split}$(a) = (a_h, a_l)$ with $a = a_h + a_l$
and $a_h$ having $\lfloor p/2 \rfloor$ trailing zeros is the
source of the rescaling logic peculiar to \texttt{edd\_real}; we
present it in Section~\ref{sec:split}.

\textsc{TwoSqr}$(a)$ is the specialization of \textsc{TwoProd} to
$b=a$ and requires only one \textsc{Split}.

\section{Representation and Invariants}
\label{sec:repr}

An \texttt{edd\_real} stores a two-component array of binary80
words:
\begin{lstlisting}
using edd_word = _Float64x;
struct edd_real {
  edd_word x[2];      // x[0] = hi, x[1] = lo
  ...
};
\end{lstlisting}
with the semantic value
\[
  a \;=\; a_0 + a_1, \qquad a_i := \texttt{x[$i$]}.
\]
A valid \texttt{edd\_real} is \emph{normalized}:
\[
  a_0 = \mathrm{fl}(a_0 + a_1), \qquad
  |a_1| \;\le\; \tfrac{1}{2}\,\mathrm{ulp}_{80}(a_0),
\]
that is, the two limbs are non-overlapping in binary80.  We use
the test
\[
  |a_1| \;\le\; \tfrac{1}{2} \cdot 2^{-63} \cdot |a_0|
\]
as an effective invariant check.  The constants \texttt{\_eps},
\texttt{\_ndigits}, and the numeric limits specializations follow:
\begin{lstlisting}
const edd_word edd_real::_eps = ldexpx(1.0L, -126);
const int edd_real::_ndigits = 38;
namespace std {
  template<> class numeric_limits<edd_real>
      : public numeric_limits<long double> {
    static const int digits   = 126;
    static const int digits10 = 37;
    ...
  };
}
\end{lstlisting}
The effective unit roundoff of \texttt{edd\_real} is therefore
$\varepsilon = 2^{-126}$, and worst-case representability covers
roughly 37 safely round-trippable decimal digits.

\section{The Splitter with Rescaling}
\label{sec:split}

The Dekker--Veltkamp split requires a constant
\[
  C \;=\; 2^{\lceil p/2 \rceil} + 1 \;=\; 2^{32} + 1.
\]
The naive split,
\[
  t \gets C \cdot a, \quad
  a_h \gets t - (t - a), \quad
  a_l \gets a - a_h,
\]
is exact when $C \cdot a$ is representable.  For
\texttt{edd\_real} this requires $|a| \lesssim 2^{E_{\max} - 33}$.
A plain dd-style splitter therefore overflows near the top of the
binary80 exponent range, and is especially problematic because
\textsc{TwoProd} is called from high-precision transcendental
functions that legitimately reach very large intermediate values.

The implementation guards the naive split with a magnitude test
and a temporary downscaling:

\begin{lstlisting}
inline void split(edd_word a, edd_word& hi, edd_word& lo) {
  edd_word temp;
  if (fabsx(a) > split_thresh()) {       // |a| > 2^(E_max - 33)
    a   = ldexpx(a, -33);
    temp = splitter() * a;
    hi   = temp - (temp - a);
    lo   = a - hi;
    hi   = ldexpx(hi, 33);
    lo   = ldexpx(lo, 33);
  } else {
    temp = splitter() * a;
    hi   = temp - (temp - a);
    lo   = a - hi;
  }
}
\end{lstlisting}
Here \texttt{splitter()} is $2^{32} + 1$ and
\texttt{split\_thresh()} is $2^{E_{\max} - 33}$.  Exact
\texttt{ldexpf64x} scaling by $\pm 33$ introduces no roundoff,
so the rescaled split is mathematically identical to the
unrescaled one; the guard exists only to keep the multiplication
by $C$ finite.  The exponent $33$ is chosen as the smallest
constant strictly greater than $\lceil p/2 \rceil = 32$, which is
sufficient to absorb both the factor $C$ and the subsequent
$t - a$.

\begin{remark}
The corresponding \texttt{dd\_real} splitter uses $C = 2^{27}+1$
and a threshold of $2^{E_{\max}-28}$.  The constants differ only
in the value of $\lceil p/2 \rceil$; the structure of the rescale
is identical.  The QD-codebase terminology for the scaling
exponent is \texttt{QD\_EDD\_SPLIT\_SCALE\_BITS}, defined as
$33$.
\end{remark}

\section{Renormalization}
\label{sec:renorm}

Because \texttt{edd\_real} has only two limbs, renormalization
reduces to a single \textsc{QuickTwoSum} call.  The
implementation uses
\begin{lstlisting}
inline void renorm2(edd_word& c0, edd_word& c1) {
  if (isinfx(c0)) return;
  c0 = quick_two_sum(c0, c1, c1);
}
\end{lstlisting}
which enforces $c_0 = \mathrm{fl}(c_0 + c_1)$ and leaves $c_1$ as
the corresponding error term.  A 4-component variant
\texttt{renorm4} is also provided: it is used when converting
from \texttt{qd\_real} to \texttt{edd\_real}, where we first form
a four-component sum by limb-wise truncation and then compress it
to two normalized components (Section~\ref{sec:mixed}).

\section{Addition}
\label{sec:add}

\subsection{\texttt{edd\_real $+$ edd\_word}}

Adding a scalar $b$ to $a = (a_0, a_1)$:
\begin{algorithmic}[1]
\State $(s_1, s_2) \gets \textsc{TwoSum}(a_0, b)$
\State $s_2 \gets s_2 + a_1$
\State $(s_1, s_2) \gets \textsc{QuickTwoSum}(s_1, s_2)$
\State \Return $(s_1, s_2)$
\end{algorithmic}
Cost: 1 \textsc{TwoSum} ($6$ flops) + 1 addition + 1
\textsc{QuickTwoSum} ($3$ flops) $=10$ flops.

\subsection{\texttt{edd\_real $+$ edd\_real}}

We use the IEEE-style (``sloppy'') sum that the QD library adopts
for \texttt{dd\_real}:
\begin{algorithmic}[1]
\State $(s_1, s_2) \gets \textsc{TwoSum}(a_0, b_0)$
\State $(t_1, t_2) \gets \textsc{TwoSum}(a_1, b_1)$
\State $s_2 \gets s_2 + t_1$
\State $(s_1, s_2) \gets \textsc{QuickTwoSum}(s_1, s_2)$
\State $s_2 \gets s_2 + t_2$
\State $(s_1, s_2) \gets \textsc{QuickTwoSum}(s_1, s_2)$
\State \Return $(s_1, s_2)$
\end{algorithmic}
Cost: 2 \textsc{TwoSum} + 2 \textsc{QuickTwoSum} $=18$ flops plus
two scalar additions.  This is the same arithmetic pattern as the
QD library's \texttt{sloppy\_add}.

Subtraction \texttt{edd\_real $-$ edd\_real} uses the identical
structure with \textsc{TwoDiff} in place of \textsc{TwoSum}.

\section{Multiplication}
\label{sec:mul}

\subsection{\texttt{edd\_real $\times$ edd\_word}}

\begin{algorithmic}[1]
\State $(p_1, p_2) \gets \textsc{TwoProd}(a_0, b)$
\State $p_2 \gets p_2 + a_1 b$
\State $(p_1, p_2) \gets \textsc{QuickTwoSum}(p_1, p_2)$
\State \Return $(p_1, p_2)$
\end{algorithmic}
Cost: 1 \textsc{TwoProd} + 1 mul-add + 1 \textsc{QuickTwoSum}.

\subsection{\texttt{edd\_real $\times$ edd\_real}}

\begin{algorithmic}[1]
\State $(p_1, p_2) \gets \textsc{TwoProd}(a_0, b_0)$
\State $p_2 \gets p_2 + a_0 b_1 + a_1 b_0$
  \Comment{cross-terms; $a_1 b_1$ is beyond precision}
\State $(p_1, p_2) \gets \textsc{QuickTwoSum}(p_1, p_2)$
\State \Return $(p_1, p_2)$
\end{algorithmic}
Cost: 1 \textsc{TwoProd} + 2 multiplications + 2 additions + 1
\textsc{QuickTwoSum}.

The $a_1 b_1$ product is at most
$\mathrm{ulp}_{80}(a_0) \cdot \mathrm{ulp}_{80}(b_0)$ in magnitude
and therefore lies below the rounding threshold of $p_1$.  It is
dropped, which is the source of the $\approx 2$-bit loss relative
to the ideal $2p = 128$-bit product; the stored accuracy is
$2p - 2 = 126$ bits, consistent with
$\texttt{\_eps} = 2^{-126}$.

\section{Division}
\label{sec:div}

Division uses the classical three-digit long-division scheme, with
an additional rescaling pattern analogous to (but tuned
differently from) the splitter.  The rescaling is necessary
because the quotient $q_1 = a_0 / b$ can overflow when $|a_0|$
approaches the binary80 maximum: dividing by a small $b$ amplifies
the magnitude of $a_0$.  The threshold used is
\[
  T_{\mathrm{div}} \;=\; 2^{E_{\max} - \texttt{\_\_FLT64X\_MANT\_DIG\_\_}}
  \;=\; 2^{16384 - 64} \;=\; 2^{16320},
\]
slightly below the overflow boundary by exactly one mantissa
width.

The implementation of \texttt{edd\_real / edd\_real} is

\begin{lstlisting}
inline edd_real operator/(const edd_real& a, const edd_real& b) {
  const bool rescale = edd_real_div_needs_rescale(a.x[0]);
  const edd_real aa  = rescale
                     ? mul_pwr2(a, ldexpx(1.0L, -64))
                     : a;
  edd_word q1 = aa.x[0] / b.x[0];
  edd_real r  = aa - q1 * b;
  edd_word q2 = r.x[0] / b.x[0];
  r          -= q2 * b;
  edd_word q3 = r.x[0] / b.x[0];
  q1 = quick_two_sum(q1, q2, q2);
  r  = edd_real(q1, q2) + q3;
  return rescale ? mul_pwr2(r, ldexpx(1.0L, 64)) : r;
}
\end{lstlisting}

The structure mirrors the QD library's \texttt{sloppy\_div} for
\texttt{dd\_real}: three trial quotients are accumulated, and the
first two are combined by a \textsc{QuickTwoSum} before the final
correction is added.  Because the rescale is by an exact power of
two, it cancels when the rescale flag is true on entry and
restored on exit.

\section{Square Root}
\label{sec:sqrt}

The square root uses a single Karp-style refinement step
\cite{Karp1997}.  Given $a > 0$, seed $x = 1/\sqrt{a_0}$ in
binary80 and refine once:
\[
  s \;=\; a_0 x, \qquad
  \sqrt{a} \;\approx\; s \;+\; (a - s^2)_0 \cdot (x/2).
\]
Reorganized in error-free form,
\begin{lstlisting}
inline edd_real sqrt_karp(const edd_real& a) {
  edd_word x  = 1.0L / sqrtx(a[0]);
  edd_word ax = a[0] * x;
  return edd_real::add(
      ax,
      (a - edd_real::sqr(ax))[0] * (x * 0.5L));
}
\end{lstlisting}
The outer call handles signs, zero, and an overflow guard:

\begin{lstlisting}
edd_real sqrt(const edd_real& a) {
  if (a.is_zero())     return 0.0L;
  if (a.is_negative()) { error(...); return _nan; }
  if (fabsx(a[0]) > sqrt_rescale_thresh()) {
    // |a[0]| > 2^(E_max - 3)
    return mul_pwr2(sqrt_karp(mul_pwr2(a, 0.25L)), 2.0L);
  }
  return sqrt_karp(a);
}
\end{lstlisting}
The threshold $2^{E_{\max} - 3}$ is necessary because the
refinement step forms $s^2$: for $|a|$ near the binary80 maximum,
the product $a_0 \cdot x$ is in the square-root regime
$2^{E_{\max}/2}$ and $s^2$ would overflow if $x$ is not quite
accurate.  Scaling $a$ down by $4$ (one extra bit of exponent
margin on each side of the square root) avoids the issue, and the
final result is scaled back by $2$.

Convergence of one Karp step to full binary80 precision is
quadratic, so the seed accuracy of $\approx 63$ bits from the
native \texttt{sqrtf64x} is amplified to $\approx 126$ bits in a
single refinement.  The hardware \texttt{sqrtf64x} is
correctly-rounded on every supported platform, which is what
makes a single-step refinement sufficient.

\section{Transcendental Functions}
\label{sec:trans}

The transcendental implementations follow a \emph{hybrid} design.
For each function, the code first tests whether the argument is
safely within the \texttt{qd\_real} range:
\begin{lstlisting}
inline bool qd_safe_argument(const edd_real& a) {
  return a.isfinite() && abs(a) <= to_edd_real(qd_real::_safe_max);
}
\end{lstlisting}
If so, it converts $a$ to \texttt{qd\_real}, delegates to the
high-quality \texttt{qd\_real} implementation, and converts back.
This leverages QD's existing, well-tested range reduction and
Taylor-series machinery without re-deriving it for binary80 limbs.
For arguments outside the \texttt{qd\_real} safe range, a
\texttt{edd\_real}-native scheme takes over.

\subsection{\texttt{exp}}

For any $a$, the native implementation uses the classical
argument reduction
\[
  a \;=\; m \ln 2 + k \cdot r, \qquad
  k = 2^{-12}, \quad |r| \le \tfrac{k}{2}\ln 2,
\]
and computes $\mathrm{expm1}(r)$ by Taylor series,
then folds in $k$ by repeated squaring and in $m$ by
\texttt{ldexp}:
\begin{algorithmic}[1]
\State $m \gets \lfloor a_0 / (\ln 2)_0 + 0.5 \rfloor$
\State $r \gets (a - m \ln 2) / 4096$
\State $s \gets \textsc{expm1Taylor}(r)$
\For{$i = 1$ to $12$}
  \State $s \gets 2s + s^2$ \Comment{$s$ represents $e^r - 1$}
\EndFor
\State \Return $\textsc{ldexp}(s + 1, m)$
\end{algorithmic}
Overflow and underflow are pre-tested against $a_0 \gtrless
c \ln 2$ with $c = 16384$ for overflow and $c = -16445$ for
underflow; these bounds come directly from the binary80 exponent
range.

\subsection{\texttt{log}}

A Newton iteration of the equation $\exp(x) - a = 0$ converges
cubically for correct seeds, so three iterations suffice from
the binary80 \verb|__builtin_logf64x| seed:
\begin{algorithmic}[1]
\State $(m, e) \gets \texttt{frexpf64x}(a_0)$
\State $x \gets \ln m + e \ln 2$
\For{$i = 1$ to $3$}
  \State $x \gets x + a \cdot e^{-x} - 1$
\EndFor
\State \Return $x$
\end{algorithmic}
The update $x \gets x + a e^{-x} - 1$ is derived from
$f(x) = e^x - a$, $f'(x) = e^x$, so $x_{\mathrm{new}} = x -
(e^x - a)/e^x = x + ae^{-x} - 1$.  The $\log_{10}$ function is
implemented as $\log(a) / \log 10$.

\subsection{\texttt{sincos}}

For $|a| \le \texttt{qd\_real::\_safe\_max}$, the implementation
converts to \texttt{qd\_real} and delegates.  For arguments
outside this range, a native routine performs a two-stage
reduction modulo $2\pi$ and $\pi/2$ followed by a single
$\pi/16$-table lookup:

\begin{algorithmic}[1]
\State $z \gets \mathrm{nint}(a / 2\pi)$
\State $r \gets a - 2\pi z$
\State $q \gets \lfloor r_0 / (\pi/2)_0 + 0.5 \rfloor$
\State $t \gets r - (\pi/2) q$,\ $j \gets q$
\State $q \gets \lfloor t_0 / (\pi/16)_0 + 0.5 \rfloor$
\State $t \gets t - (\pi/16) q$,\ $k \gets q$
\State $(\sin t, \cos t) \gets \textsc{SinCosTaylor}(t)$
\State combine with $\sin(k\pi/16)$, $\cos(k\pi/16)$ from table
\State rotate by $j\pi/2$
\end{algorithmic}
The tables are held as pre-normalized \texttt{qd\_real} constants
and converted to \texttt{edd\_real} on use.  Only four
$k$-positions are needed ($k = 1,2,3,4$) because symmetry handles
the remaining angles.  This matches the pattern used in the QD
library's \texttt{sincos\_qd}, adapted to the two-limb expansion.

\subsection{Inverse trigonometric and hyperbolic functions}

The inverse trigonometric functions \texttt{asin}, \texttt{acos},
and \texttt{atan} always delegate to \texttt{qd\_real} via
conversion, because their arguments are bounded by $1$ and
therefore always lie in the safe range.  The hyperbolic functions
split on magnitude: for $|a| > 0.05$, \texttt{sinh} and
\texttt{cosh} are implemented via $e^a \pm e^{-a}$ using the
native \texttt{exp}; for $|a| \le 0.05$, a Taylor series avoids
the catastrophic cancellation in $(e^a - e^{-a})/2$.
\texttt{atan2} uses a Newton iteration driven by the binary80
\verb|__builtin_atan2f64x| seed.

\section{Type Conversion and Mixed-Precision Arithmetic}
\label{sec:mixed}

Interoperability with \texttt{dd\_real} and \texttt{qd\_real} is
essential because the transcendental fallback depends on it.

\subsection{\texttt{dd\_real} $\leftrightarrow$ \texttt{edd\_real}}

\texttt{dd\_real} has two binary64 limbs whose combined precision
($\approx 106$ bits) is narrower than a single binary80 word.
The forward conversion is therefore exact, given a renormalization:
\begin{lstlisting}
edd_real::edd_real(const dd_real& dd) {
  x[0] = (edd_word) dd._hi();
  x[1] = (edd_word) dd._lo();
  renorm2(x[0], x[1]);
}
\end{lstlisting}
The reverse direction is lossy: a single binary80 word already
carries more bits than two binary64 words combined, so we
decompose each limb into a head/tail pair and ship the top four
fragments to \texttt{qd\_real} before taking its two high
components.

\subsection{\texttt{qd\_real} $\leftrightarrow$ \texttt{edd\_real}}

For the forward direction \texttt{qd\_real $\to$ edd\_real}, we
truncate from four binary64 limbs to two binary80 limbs:
\begin{lstlisting}
inline edd_real from_qd_truncate(const qd_real& a) {
  edd_word c0 = (edd_word) a[0], c1 = (edd_word) a[1];
  edd_word c2 = (edd_word) a[2], c3 = (edd_word) a[3];
  renorm4(c0, c1, c2, c3);
  return edd_real(c0, c1);
}
\end{lstlisting}
The \texttt{renorm4} cascade compresses the four limbs into two
non-overlapping components while discarding bits that fall below
the binary80 rounding threshold.

For the reverse direction \texttt{edd\_real $\to$ qd\_real}, the
routine splits each \texttt{edd\_word} into a binary64 head and a
binary64 tail by conversion arithmetic:
\begin{lstlisting}
inline void word_to_doubles(edd_word x, double& hi, double& lo) {
  hi = static_cast<double>(x);
  lo = static_cast<double>(x - (edd_word) hi);
}
\end{lstlisting}
The resulting four binary64 numbers are handed to
\texttt{qd\_real::renorm}.  Because binary80 rounding of the
subtraction \texttt{x - (edd\_word) hi} is exact (the head has at
most 53 significant bits, so the tail fits within one binary80
mantissa), the decomposition is error-free.

\subsection{Mixed-precision operators}

All arithmetic operators accept a mixed \texttt{edd\_word}/%
\texttt{edd\_real} pair and forward to the appropriate kernel.
Comparison operators are implemented via
\begin{lstlisting}
inline bool comparison_lt(const A& a, const B& b) {
  edd_real ea(a), eb(b);
  if (ea.isnan() || eb.isnan()) return false;
  if (ea.isinf() || eb.isinf()) return ea[0] <  eb[0];
  return (ea - eb).is_negative();
}
\end{lstlisting}
This promotes \texttt{edd\_word} operands to \texttt{edd\_real}
and delegates to the sign test of the difference.  Bit-level
equality (including the $(1 - 2^{-63}, 2^{-63})$ pair with
$(1, 0)$) is correctly recognized as equal by this scheme because
\texttt{renorm2} collapses the difference to an exact zero.

\section{Constants}
\label{sec:const}

Elementary \texttt{edd\_real} constants are \emph{not} recomputed
at runtime.  Instead, the well-established \texttt{qd\_real}
reference values are truncated through \texttt{renorm4}:
\begin{lstlisting}
const edd_real edd_real::_pi    = to_edd_real(qd_real::_pi);
const edd_real edd_real::_2pi   = to_edd_real(qd_real::_2pi);
const edd_real edd_real::_pi2   = to_edd_real(qd_real::_pi2);
const edd_real edd_real::_pi4   = to_edd_real(qd_real::_pi4);
const edd_real edd_real::_3pi4  = to_edd_real(qd_real::_3pi4);
const edd_real edd_real::_e     = to_edd_real(qd_real::_e);
const edd_real edd_real::_log2  = to_edd_real(qd_real::_log2);
const edd_real edd_real::_log10 = to_edd_real(qd_real::_log10);
\end{lstlisting}
The \texttt{qd\_real} constants carry more than 200 bits of
accuracy, comfortably exceeding the 126 bits of
\texttt{edd\_real}; truncation therefore loses no
representable bits.  Independent computation of these constants
(e.g.\ by AGM as in the example \texttt{examples/edd\_agm.cpp})
agrees with the truncated values to within a small multiple of
\texttt{edd\_real::\_eps}.

Special values are generated directly from binary80 primitives:
\texttt{\_nan} via \verb|__builtin_nanf64x("")|, \texttt{\_inf}
via \verb|__builtin_huge_valf64x()|.

\section{Build-System Detection}
\label{sec:build}

\texttt{edd\_real} is an optional component of the QD library; it
is compiled only when the host compiler advertises binary80
\texttt{\_Float64x} support.  The CMake probe
compiles a small program that asserts
\begin{itemize}
  \item \verb|__GNUC__| is defined,
  \item the target is x86 or x86\_64,
  \item \verb|__FLT64X_MANT_DIG__ == 64|,
  \item \verb|__FLT64X_MAX_EXP__ == 16384|, and
  \item the relevant \texttt{\_Float64x} builtins link successfully.
\end{itemize}
On success, the probe defines
\begin{center}
\begin{tabular}{ll}
\toprule
Macro & Value \\
\midrule
\texttt{QD\_HAVE\_EDD\_REAL}          & $1$ \\
\texttt{QD\_EDD\_FLT64X\_MANT\_DIG}   & $64$ \\
\texttt{QD\_EDD\_FLT64X\_MAX\_EXP}    & $16384$ \\
\texttt{QD\_EDD\_SPLIT\_BITS}         & $32$ \\
\texttt{QD\_EDD\_SPLIT\_SCALE\_BITS}  & $33$ \\
\bottomrule
\end{tabular}
\end{center}
and the Automake conditional \texttt{HAVE\_EDD\_REAL} becomes
true.  The Makefile then adds \texttt{edd\_real.cpp},
\texttt{edd\_const.cpp}, \texttt{edd\_trans.cpp}, and
\texttt{c\_edd.cpp} to \texttt{libqd}.  A summary line
\texttt{edd\_real = enabled/disabled} is printed at the end of
\texttt{configure}.

On non-GNU compilers or on non-x86 targets the probe fails
silently, the macro \texttt{QD\_HAVE\_EDD\_REAL} remains
undefined, and the binary64-based \texttt{dd\_real},
\texttt{td\_real}, and \texttt{qd\_real} types continue to be
built as usual.  A header guard in \texttt{edd\_real.h} converts
any accidental inclusion into a compile-time error.

\section{C Bindings and the Test Suite}
\label{sec:test}

\subsection{C wrapper}

A thin C interface is provided by \texttt{include/qd/c\_edd.h}.
The wrapper functions take \texttt{const \_Float64x *} input and
\texttt{\_Float64x *} output and internally reinterpret them as
\texttt{edd\_real} values.  All of the arithmetic, comparison,
transcendental, and string-I/O operators are exposed, together
with the constants \texttt{c\_edd\_pi}, \texttt{c\_edd\_2pi}, and
\texttt{c\_edd\_epsilon}.  The smoke test in \texttt{tests/c\_test.c}
verifies $\sin^2 x + \cos^2 x = 1$, $\log(\exp x) = x$, and
$|\tanh x| < 1$ through the C interface.

\subsection{C++ test suite}

The C++ test suite \texttt{tests/qd\_test.cpp} adds an
\texttt{-edd} mode alongside the existing
\texttt{-dd}/\texttt{-td}/\texttt{-qd} switches.  The
\texttt{EddTestSuite} class contains 18 numbered tests covering:

\begin{center}
\begin{tabular}{rl}
\toprule
\# & Purpose \\
\midrule
1  & Normalization after $+$, $-$, $\times$, $/$, $\sqrt{\,}$ \\
2  & Cancellation-sensitive $a-b$ and round-trip \\
3  & Mixed-magnitude multiplication vs.\ \texttt{qd\_real} \\
4  & Basic division, $1/3$, inverse check \\
5  & $\sqrt{}$ on exact squares \\
6  & $\sqrt{}$ near $1$ \\
7  & String parse / format round trip \\
8  & Constructor / scalar conversion \\
9  & Constants \\
10 & String I/O with \texttt{inf}, \texttt{nan}, signed zero \\
11 & Mixed-mode \texttt{edd\_word} arithmetic \\
12 & \texttt{dd}/\texttt{qd} conversion paths \\
13 & Overflow / underflow boundaries \\
14 & \texttt{exp}/\texttt{log} across safe and wide ranges \\
15 & Trig identities, agreement with \texttt{qd\_real} oracle \\
16 & Inverse-trig consistency \\
17 & Large-argument reduction, hyperbolic identities \\
18 & Comparison normalization for overlapping inputs \\
\bottomrule
\end{tabular}
\end{center}

Tolerances are expressed as small multiples of
\texttt{edd\_real::\_eps}, with the \texttt{qd\_real} value as
reference oracle.  The test binary also exercises a stress case
in which $\sin(a)$ and $\cos(a)$ must agree with the
\texttt{qd\_real} result at $a = 10^{40}\pi + 0.125$; the native
range-reduction path is active for this argument.

Before running the \texttt{edd\_real} tests, the test harness
calls \texttt{fpu\_fix\_end} to release the x87 round-to-double
control word that the \texttt{dd}/\texttt{qd}/\texttt{td} tests
install.  This is critical: \texttt{edd\_real} needs the x87 FPU
in 80-bit mode, whereas \texttt{dd\_real} requires the 53-bit
round-to-double mode for its error-free transformations to be
provably exact.

\subsection{AGM example}

The file \texttt{examples/edd\_agm.cpp} computes $\pi$, $\log 2$,
$\log 10$, $\log_{10} 2$, and $\log_{10} e$ from scratch by the
Brent--Salamin AGM iteration and the Brent--McMillan
logarithm-via-AGM formula.  The program compares the AGM results
to the \texttt{edd\_real} built-in constants and reports the
errors in units of \texttt{edd\_real::\_eps}.  With $8$ AGM
iterations and a scaling shift of $80$, all five constants agree
with their references to within a few times \texttt{\_eps}, i.e.\
within a few ulps of full \texttt{edd\_real} precision.

\section{Discussion and Conclusion}
\label{sec:conclusion}

We have described \texttt{edd\_real}, a two-limb extended-double-%
double floating-point type based on the GNU C++ \texttt{\_Float64x}
binary80 format.  Its unit roundoff $2^{-126}$ corresponds to
$\approx 38$ decimal digits, placing it between \texttt{dd\_real}
and \texttt{td\_real} in precision but exceeding both in dynamic
range.  The two-limb structure keeps the per-operation cost close
to \texttt{dd\_real}, while the hybrid transcendental strategy
lets the implementation reuse \texttt{qd\_real}'s extensively
tested range reduction whenever the argument is small enough.

The distinguishing technical features of the implementation are
(a) the splitter rescaling at $2^{-33}$, derived from
$\lceil p/2 \rceil + 1 = 33$ for binary80; (b) a division
rescaling at the threshold $2^{E_{\max}-64}$ keyed to the full
mantissa width; and (c) a square-root rescaling at the threshold
$2^{E_{\max}-3}$ keyed to the squaring step inside Karp's
refinement.  These three thresholds are all exactly-representable
powers of two, so the rescale/restore steps are arithmetically
invisible; they exist solely to keep intermediate products finite
near the top of the binary80 exponent range.

\texttt{edd\_real} is Phase~1 complete: it is correct, well-tested,
and interoperates cleanly with the rest of the QD library, but its
transcendental implementations largely delegate to
\texttt{qd\_real} for small arguments.  A Phase~2 effort could
develop fully native \texttt{edd\_real} implementations of the
range-reduction machinery, avoiding the round-trip through
\texttt{qd\_real}; this would unlock the full binary80 dynamic
range for sine, cosine, and arctangent at near-full
\texttt{edd\_real} precision, and is a natural continuation of the
work described here.

\begin{thebibliography}{9}

\bibitem{QD}
Y.~Hida, X.~S.~Li, and D.~H.~Bailey.
\emph{Library for Double-Double and Quad-Double Arithmetic}.
Lawrence Berkeley National Laboratory, 2000--2023.
\url{https://github.com/BL-highprecision/QD}.

\bibitem{HidaLiBailey2000}
Y.~Hida, X.~S.~Li, and D.~H.~Bailey.
Algorithms for quad-double precision floating-point arithmetic.
In \emph{Proc.\ 15th IEEE Symposium on Computer Arithmetic},
pp.\ 155--162, 2001.

\bibitem{Dekker1971}
T.~J.~Dekker.
A floating-point technique for extending the available precision.
\emph{Numerische Mathematik}, 18(3):224--242, 1971.

\bibitem{Knuth1997}
D.~E.~Knuth.
\emph{The Art of Computer Programming, Volume 2: Seminumerical
Algorithms}, 3rd ed.
Addison-Wesley, 1997.

\bibitem{Karp1997}
A.~H.~Karp and P.~Markstein.
High-precision division and square root.
\emph{ACM Trans.\ Math.\ Softw.}, 23(4):561--589, 1997.

\bibitem{Intel80}
Intel Corporation.
\emph{Intel 64 and IA-32 Architectures Software Developer's
Manual, Vol.\ 1: Basic Architecture}.  Chapter on the x87 FPU.

\bibitem{BrentSalamin}
R.~P.~Brent.
Multiple-precision zero-finding methods and the complexity of
elementary function evaluation.
In \emph{Analytic Computational Complexity}, pp.\ 151--176,
Academic Press, 1975.

\end{thebibliography}

\end{document}
